% ==================== 3.6 本章小结 ====================
\section{本章小结}
\label{sec:ch3_summary}

本章针对月面环境的特殊性和复杂性，系统研究了月面地形、月壤环境、光照与温度环境的数字孪生建模方法。主要工作和结论如下：

（1）在月面地形数字孪生建模方面，基于LOLA等遥感数据构建了三维地形模型，采用规则格网和不规则三角网两种表达方式，实现了坡度、坡向、粗糙度等地形特征的自动提取，以及撞击坑、岩石等障碍物的识别。

（2）在月壤环境数字孪生建模方面，分析了月壤的颗粒组成、形状特征、密度分布等物理特性，建立了基于Drucker-Prager准则的月壤弹塑性本构模型和基于离散元的颗粒力学模型，实现了月壤参数空间分布的建模与预测。

（3）在光照与温度环境建模方面，分析了月球特有的太阳运动规律，建立了基于光线追踪的光照场模型和基于热平衡方程的温度场模型，实现了月面任意位置、任意时刻光照和温度状态的精确计算。

（4）通过实验验证了各模型的有效性。地形模型能够准确再现月面地形特征，月壤模型预测的车轮沉陷与实验数据误差在15\%以内，温度模型预测值与观测数据误差在10\%以内。

本章建立的月面环境数字孪生模型为巡视器动力学建模、路径规划、避障决策等后续研究奠定了环境基础。
